\documentclass[11pt]{article}
\usepackage[a4paper,margin=0.8in]{geometry}
\usepackage[T1]{fontenc}
\usepackage{tgtermes}
\usepackage{enumitem}
\usepackage[hidelinks]{hyperref}
\usepackage{microtype}
\pagestyle{empty}

% Add these lines to fix underfull hbox warnings
\tolerance=1
\emergencystretch=\maxdimen
\hyphenpenalty=10000
\hbadness=10000

\begin{document}

\begin{center}
  {\LARGE \textbf{ YESHWANTH REDDY GURREDDY }}\\[4pt]
  Robotics / Embedded Controls Engineer \quad $\vert$ \quad Tempe, AZ\\[2pt]
  \href{mailto:yeshwanth.rg060694@gmail.com}{ yeshwanth.rg060694@gmail.com } \quad $\vert$ \quad \href{https://github.com/https://github.com/YeshRobo}{github/https://github.com/YeshRobo}
\end{center}

\vspace{6pt}
\noindent\hrulefill

\vspace{6pt}
\section*{Summary}
Experienced robotics engineer pursuing an M.S. in Robotics and Autonomous Systems at Arizona State University (May 2026). Background in UAV gimbals, embedded motor control (SVPWM, FOC), and embedded system architecture. Strong hands-on experience with microcontrollers, hardware debugging, and cross-functional engineering delivery.

\section*{Experience}

\textbf{ Master's Thesis Research (in progress) } \hfill 2024-08 -- present\\
Arizona State University \\

\begin{itemize}[leftmargin=*,noitemsep]
  
  \item Conducting thesis research on effort-aware adaptive control for human–exoskeleton interaction, focusing on separating motion-driven effort from load-support effort.
  
  \item Reviewing and comparing baseline modeling and residual methods for estimating load-related torque/effort from EMG and kinematic inputs.
  
  \item Studying safety-constrained control formulations (e.g., CLF/CBF/QP-based methods) as candidate approaches for enforcing stability and safety constraints during user motion.
  
  \item Developing software prototypes and simulation-based checks to validate modeling/control assumptions and define an evaluation plan (in progress).
  
\end{itemize}

\vspace{4pt}

\textbf{ Graduate Research Assistant } \hfill 2025-07 -- 2025-12\\
Arizona State University \\

\begin{itemize}[leftmargin=*,noitemsep]
  
  \item Ran shoulder exoskeleton reaching/tracking experiments on 10 human subjects, executing sessions and logging joint position/velocity, EMG, and human–robot interaction torque (6-axis F/T) at 200 Hz (~30 minute sessions) using Python.
  
  \item Expanded the study from a pilot dataset to a 10-subject dataset to strengthen model-vs-baseline comparison and support follow-on paper revision.
  
  \item Studied safety-constrained control formulations (e.g., CBF-QP safety filters) and ensured collected signals (interaction torque/kinematics) were suitable for downstream constraint/safety analysis.
  
  \item Identified that the original visual feedback did not clearly differentiate model vs baseline; redesigned target presentation (patterns/spatial layouts) and implemented the revised feedback logic in Python within the experimental workflow.
  
\end{itemize}

\vspace{4pt}

\textbf{ Design Engineer } \hfill 2023-05 -- 2024-07\\
Future Electronics \\

\begin{itemize}[leftmargin=*,noitemsep]
  
  \item Developed a proof-of-concept quadcopter flight controller using the CY8CKIT-062S2-43012 and CY8CKIT-028-SENSE development kits to demonstrate feasibility for small UAV control on a PSoC-based platform.
  
  \item Supported customers using Infineon and STMicroelectronics MCUs/MPUs by analyzing technical queries and providing practical guidance on hardware–firmware integration and debugging.
  
  \item Expanded embedded experience across multiple MCU ecosystems and evaluation kits, complementing prior product-focused controls and firmware work.
  
\end{itemize}

\vspace{4pt}

\textbf{ Engineer (PTC ThingWorx) } \hfill 2021-05 -- 2023-04\\
DELOPT \\

\begin{itemize}[leftmargin=*,noitemsep]
  
  \item Led adoption of PTC ThingWorx as a newly introduced platform, learning it from scratch and becoming a key contributor to its internal implementation.
  
  \item Migrated a legacy retail monitoring application to PTC ThingWorx end-to-end, delivering equivalent functionality while enabling better scalability and modernization.
  
  \item Integrated ThingWorx with SQL databases to ingest and manage time-series data from retail assets across ~10 sites and ~70 assets (prototype scale).
  
  \item Developed backend ThingWorx Services and SQL queries to support real-time and near-real-time dashboards with periodic refresh (~10-minute intervals) to keep data current.
  
  \item Designed and implemented data models using Things, ThingShapes, ThingTemplates, Properties, Services, and InfoTables to support reusable and scalable dashboard services.
  
  \item Built reusable dashboard services and components to standardize data access patterns across multiple retail sites and assets.
  
  \item Implemented access control, user/group management, and environment configurations; supported deployments and troubleshooting across dev, test, and prod environments.
  
  \item Mentored junior engineers on ThingWorx concepts, data modelling, and dashboard/service development, becoming the internal go-to person for the platform.
  
  \item Collaborated closely with controls, QA, and product teams using verbal requirements, iterative demos, and feedback cycles to refine features.
  
\end{itemize}

\vspace{4pt}

\textbf{ Senior Engineer } \hfill 2019-05 -- 2021-04\\
DELOPT \\

\begin{itemize}[leftmargin=*,noitemsep]
  
  \item Owned the system-level design of a 2-axis stabilized UAV EO gimbal, translating requirements into a complete architecture (control blocks, sensing chain, actuation chain, compute, power, and comms).
  
  \item Created high-level block diagrams and signal flow definitions (what runs where, what sensor feeds which loop, what updates at what rate), aligning mechanical/electrical/firmware teams on a single integration plan.
  
  \item Drove placement decisions for motors, IMU, and encoders to support stability and sensing quality (axis alignment, reference frames, vibration exposure, cable strain, and serviceability constraints).
  
  \item Worked with mechanical design to ensure sensor mounting orientation and axis mapping were consistent with control assumptions, reducing rework during integration and tuning.
  
  \item Defined the wiring and interconnect strategy: connector selection approach, pinouts, harness routing guidance, and separation of noisy motor/power lines from sensitive IMU/encoder signals.
  
  \item Produced/maintained wiring diagrams and interface notes used during prototype builds and troubleshooting (what connects where, expected signals, and sanity-check steps).
  
  \item Led and supported hardware bring-up from first power-on through closed-loop control: continuity checks, motor phase/rotation verification, encoder direction validation, IMU sanity checks, and safe commissioning sequence.
  
  \item Acted as the cross-disciplinary debugger during integration—diagnosing issues spanning wiring, sensor mounting/orientation, electrical noise, firmware timing, and control-loop stability.
  
  \item Used bench instrumentation and logging (e.g., scope/serial telemetry-style workflows) to isolate root causes like oscillations, drift, jitter/noise on sensor lines, sign/axis mismatches, and intermittent encoder/IMU faults.
  
  \item Coordinated fixes with EE/ME teams (rewire/reroute, mounting changes, shielding/grounding tweaks, firmware updates) and verified the fix through repeatable tests.
  
  \item Designed and implemented embedded C firmware on a TI TMS320F28379 using Code Composer Studio, including PWM generation, sensor drivers, control task scheduling, and real-time safety handling.
  
  \item Implemented motor-control and stabilization algorithms from scratch (SVPWM, FOC, cascaded current/stabilization/position loops), including saturation handling and practical tuning hooks for hardware iteration.
  
  \item Added gimbal-specific protections and behaviors (gimbal-lock avoidance logic, soft limits / safe operating logic where applicable) to keep the system stable and operable across edge cases.
  
  \item Communicated design tradeoffs, integration risks, and test findings to project managers and engineering leadership, keeping execution aligned across fast-moving prototype cycles.
  
\end{itemize}

\vspace{4pt}

\textbf{ Engineer } \hfill 2017-01 -- 2019-05\\
DELOPT \\

\begin{itemize}[leftmargin=*,noitemsep]
  
  \item Served as the primary engineer for integrating a 2-axis electro-optics gimbal payload on a UAV platform, working with custom interface hardware, a video tracker, and the Basecam gimbal controller.
  
  \item Developed embedded C firmware for a microcontroller-based interface board, enabling reliable communication between the payload, video tracker, and gimbal controller (tested and debugged using serial terminals).
  
  \item Configured and tuned PID control loops for the 2-axis gimbal using the Basecam GUI to improve pointing stability and reduce overshoot in the tracking loop.
  
  \item Conducted office-based tracking tests with indoor targets and supported a design intended for medium-range (~200 m) electro-optical target tracking on UAVs.
  
\end{itemize}

\vspace{4pt}


\section*{Education}

\textbf{ M.S. in Robotics and Autonomous Systems } --- Arizona State University \hfill 2024-08--2026-05\\

\textbf{ B.Tech in Electronics and Communication Engineering } --- Amrita School of Engineering \hfill 2012-05--2016-06\\




\section*{ Professional Projects }

\textbf{ Legacy Application Migration to PTC ThingWorx } --- Bengaluru, India\\

\begin{itemize}[leftmargin=*,noitemsep]
  
  \item Led adoption of PTC ThingWorx (newly introduced internally), learning the platform from scratch and driving implementation across the application.
  
  \item Converted the legacy application end-to-end to ThingWorx, delivering parity with the legacy system and enabling future enhancements.
  
  \item Built dashboards and workflows in ThingWorx to support core user journeys and stakeholder reporting.
  
\end{itemize}


\vspace{2pt}

\textbf{ 2-Axis Stabilized UAV Gimbal } --- Bengaluru, India\\

\begin{itemize}[leftmargin=*,noitemsep]
  
  \item Led end-to-end design and development of a 2-axis UAV gimbal: component selection (gyros, position sensors, BLDC motors), embedded software architecture, and control-loop tuning.
  
  \item Implemented motor driving and stabilization control loops for accurate pointing and disturbance rejection.
  
\end{itemize}


	extbf{Impact:} weight g: 440, stabilization accuracy mrad: 1, position accuracy mrad: 1, azimuth rotation deg: 360, elevation range deg: 100\\

\vspace{2pt}


\section*{ Research Projects }

\textbf{ Master's Thesis (in progress): Effort-Aware Adaptive Control for Human–Exoskeleton Interaction } --- Tempe, AZ\\

\begin{itemize}[leftmargin=*,noitemsep]
  
  \item Investigating effort-aware adaptive control methods for human–exoskeleton interaction, emphasizing separation of motion-driven vs load-support effort.
  
  \item Analyzing baseline/residual modeling approaches for EMG-informed torque/effort estimation using kinematics as context (in progress).
  
  \item Exploring CLF/CBF/QP-style safety-constrained control as a framework for safe assistance policies (in progress).
  
\end{itemize}


\vspace{2pt}

\textbf{ Visual Feedback System for Human Movement Analysis (Shoulder Exoskeleton Study) } --- Tempe, AZ\\

\begin{itemize}[leftmargin=*,noitemsep]
  
  \item Ran shoulder exoskeleton reaching/tracking experiments on 10 human subjects, handling setup, session execution, and Python-based data logging.
  
  \item Collected synchronized joint position/velocity, EMG, and human–robot interaction torque using a 6-axis force/torque sensor at 200 Hz (~30-minute sessions) to support proposed model-vs-baseline comparison.
  
  \item Identified that the existing visual feedback did not clearly separate baseline vs proposed model behavior; redesigned target patterns and spatial layouts to elicit more discriminative movements.
  
  \item Implemented the updated visual feedback logic in Python and integrated it into the existing experimental GUI workflow.
  
\end{itemize}


\vspace{2pt}


\section*{ Academic Projects }

\textbf{ A Hybrid Framework for Real-Time and Drift-Reduced Camera Pose Estimation — Academic Project } --- Arizona State University, Tempe, AZ\\

\begin{itemize}[leftmargin=*,noitemsep]
  
  \item Analyzed a hybrid camera pose estimation framework that combines semantic visual odometry (OCC-VO) with NeRF-based pose refinement to reduce long-term drift in 6-DoF trajectories on urban driving scenes.
  
  \item Set up a prototype pipeline using existing OCC-VO and UC-NeRF components on nuScenes multi-camera sequences to study how offline refinement can improve online visual odometry estimates.
  
  \item Explored using TPV-Former–style semantic occupancy representations and feature-matching–based optimization (SuperPoint/SuperGlue with Ceres-style bundle adjustment) for spatiotemporal pose refinement.
  
  \item Evaluated the approach with standard VO metrics (ATE, RPE, reprojection error) and qualitative trajectory overlays to characterize drift behavior and typical failure modes.
  
  \item Summarized trade-offs between real-time performance, accuracy, and implementation complexity for deploying such a hybrid VO–NeRF design in autonomous-driving-type setups.
  
\end{itemize}


\vspace{2pt}

\textbf{ Exploration with Partially Observable Rewards — Comparative Study } --- Arizona State University, Tempe, AZ\\

\begin{itemize}[leftmargin=*,noitemsep]
  
  \item Performed a comparative study of two published exploration strategies for reinforcement learning with partially observable rewards: a successor-representation-based 'Beyond Optimism' method and a minimax reward-agnostic exploration approach.
  
  \item Carefully examined the formal Monitored MDP setup and documented the assumptions each method makes about reward observability, monitoring, and environment structure.
  
  \item Analyzed reported experimental results on gridworld and chain-style tasks to compare sample efficiency, robustness, and behavior under different reward-monitoring regimes.
  
  \item Constructed small illustrative tabular environments to reason about how each algorithm distributes exploration effort and which states are prioritized when feedback is partial or delayed.
  
  \item Wrote a critical report highlighting practical scenarios where SR-guided, task-specific exploration is advantageous versus settings where conservative reward-agnostic coverage is safer.
  
\end{itemize}


\vspace{2pt}

\textbf{ Safe Robot Navigation with Certificated Actor–Critic — Course Project } --- Arizona State University, Tempe, AZ\\

\begin{itemize}[leftmargin=*,noitemsep]
  
  \item Implemented a Certificated Actor–Critic (CAC) safe reinforcement learning framework that learns a Control Barrier Function–based safety critic alongside a navigation policy for a simulated ground vehicle.
  
  \item Developed a kinematic bicycle–model environment with Ackermann steering and LiDAR-style range sensing to capture realistic vehicle constraints compared to point-mass baselines.
  
  \item Trained the CAC agent in scenarios with static obstacles and moving agents and compared goal-reaching success and collision rates against a reward-shaped RL baseline.
  
  \item Ran ablation-style experiments on safety pre-training duration and policy gradient restriction to study their impact on safety guarantees and overall task performance.
  
  \item Organized the project code and documentation so that instructors and classmates could reproduce the training setup and evaluation procedures.
  
\end{itemize}


\vspace{2pt}

\textbf{ Analysis of a Bio-Inspired Myriapod Robot with Modified Spine and Jansen Legs — Foldable Robotics } --- Arizona State University, Tempe, AZ\\

\begin{itemize}[leftmargin=*,noitemsep]
  
  \item Co-designed and modeled a myriapod-inspired quadruped robot with Jansen linkage legs and a compliant spine to study how spinal motion influences leg kinematics and body stability.
  
  \item Built a MuJoCo dynamic model using CAD-derived link geometry and motor actuation, sweeping spine parameters such as length and stiffness to observe effects on yaw oscillation and forward progression.
  
  \item Fabricated a laminate prototype using Smart Composite Microstructures (SCM) techniques, generating DXF cut files and assembling multi-layer flexure joints and interchangeable spine modules.
  
  \item Conducted bench experiments with overhead video tracking to compare simulated and physical locomotion, finding that the prototype produced negligible net forward motion despite correct gait cycling.
  
  \item Analyzed likely causes of the sim-to-real gap—including foot friction, leg compliance, and spine damping—and proposed design changes such as higher-friction footpads and stiffer leg laminates.
  
\end{itemize}


\vspace{2pt}




\section*{Skills}

\textbf{Languages:} C (Embedded Systems), Python, JavaScript, SQL\\
\textbf{Frameworks/Middleware:} ROS (Robot Operating System), ROS2, MATLAB/Simulink\\

\textbf{Tools:}


  

  

  

  

  

SolidWorks, Onshape, Raspberry Pi, Arduino, NVIDIA Jetson, STM32 Microcontrollers, ESP32, Texas Instruments Microcontrollers, NVIDIA Isaac Sim, CoppeliaSim (V-REP), MuJoCo, Gymnasium, TensorFlow, Visual Studio Code, Docker\\




\textbf{Soft Skills:} Stakeholder communication, Cross-functional leadership, Requirements analysis, Technical documentation, Hands-on debugging\\


\vfill
\end{document}